\chapter{Introducción}

%%%%%%%%%%%%%%%%%%%%%%%%%%%%%%%%% Motivación %%%%%%%%%%%%%%%%%%%%%%%%%%%%%%%%%%

\section{Motivación}

En la actualidad existen principalmente dos modelos de negocio en el ámbito del diseño y/o fabricación de procesadores. Por un lado estarían aquellas compañías cuyo beneficio proviene principalmente de la venta de licencias para el uso de sus diseños, como pudiera ser el caso de ARM, al tiempo que existen otras las cuales centran su actividad empresarial tanto en el diseño como en la fabricación y venta de hardware (caso de Intel), sin perjuicio de que pudieran obtener beneficios también del pago por el uso de su propiedad intelectual con fines comerciales \cite{armBusinessModel}.

Ambos modelos se basan en la siguiente premisa: que la especificación del conjunto de instrucciones que deberá ejecutar un computador constituye por sí mismo una entidad patentable. Esto obliga a todo aquel que quisiera obtener beneficio de la comercialización de productos derivados de esas especificaciones a pagar regalías al tenedor de la propiedad intelectual, lo cual obstaculiza de manera notable el surgimiento de innovaciones que pudieran aprovechar las ya existentes en materia de conjuntos de instrucciones y arquitecturas de largo recorrido. Ejemplos de estas arquitecturas son x86-64 o ARM, cuyo éxito ha quedado demostrado históricamente.

No obstante, en los últimos años del siglo XX y principios del XXI surgieron corrientes que apostaban por la creación de arquitecturas libres con proyectos como Sparc u OpenRISC, este último aún en desarrollo \cite{openrisc}. Los beneficios del hardware libre son innegables, ya que permite experimentar con arquitecturas personalizadas sin restricciones comerciales, facilitando la educación y la investigación, agilizando los procesos de creación y prueba de prototipos, y contribuyendo, en definitiva, a democratizar el acceso a la tecnología.

No sería hasta el año 2011 \cite{geneology} cuando en la Universidad de Berkeley comenzaría a gestarse el que es sin lugar a dudas el ISA abierto que mayor renombre ha conseguido tomar hasta la fecha: RISC-V. Tanto es así que gigantes tecnológicos como Google, Huawei, Microsoft o Nvidia \cite{partners} realizan importantes aportaciones monetarias de manera regular a la fundación que mantiene y desarrolla la especificación (RISC-V International).

Esta fundación recibe, además, las contribuciones de multitud de instituciones públicas. Por ejemplo, en el ámbito europeo, la EPI\footnote{European Processor Initiative} es la iniciativa mediante la cual se canalizan los esfuerzos de la academia e industria europeas por hacer realidad el proyecto de una arquitectura abierta, robusta y capaz, que sitúe a la Unión Europea a la vanguardia de la computación de altas prestaciones. Es indudable el aporte de esta iniciativa a la soberanía tecnológica y digital de Europa donde, recientemente, se ha aprobado el \textit{Chips Act 2.0}\cite{chips-act}: una propuesta en la que se recogen medidas para impulsar la industria de los microprocesadores, apoyando la producción dentro de la Unión Europea y reduciendo la dependencia de agentes externos.

En conclusión, el principal motivo que ha impulsado el desarrollo de este trabajo ha sido la posibilidad de conocer en profundidad y trabajar con una de las tecnologías clave para el futuro tecnológico de la Unión Europea.

%%%%%%%%%%%%%%%%%%%%%%%%%%%%%%%%% Objetivos %%%%%%%%%%%%%%%%%%%%%%%%%%%%%%%%%%%

\section{Objetivos}

El principal objetivo del proyecto es la implementación de un computador sencillo, mononúcleo, segmentado y en orden, que pueda ejecutar el conjunto de instrucciones definido en la especificación del ISA abierto RISC-V, en su versión de 32 bits. Para ello se partirá de un diseño ya existente, monociclo e incompleto, elaborado por el investigador de la Universidad Nacional Cheng Kung, Ching-Chun (Jim) Huang, quien lo emplea como material docente en una asignatura de arquitectura y tecnología de computadores impartida en esa misma universidad.

El diseño de partida se encuentra descrito en el lenguaje Chisel, que es un lenguaje de propósito específico o DSL (Domain-Specific Language) construido sobre el lenguaje de programación Scala, este último, un lenguaje de propósito general o GPL (General Purpose Language). Esto requerirá las adquisición, por parte del autor del presente trabajo, del conocimiento sobre ambos lenguajes que le permita, a posteriori, afrontar la tarea de completar el pipeline del mencionado diseño original para, más adelante, modificarlo de manera que este se encuentre segmentado en cinco etapas.

La segmentación del diseño original implicará, en primer lugar, un estudio profundo de la estructura interna del procesador, analizando el comportamiento y el conexionado de los diferentes bloques en los que se encuentra organizada la arquitectura del núcleo para, más adelante, integrar en él los registros de etapa y la unidad de riesgos con los que se transformará la arquitectura monociclo en una arquitectura en pipeline.

Durante todo el proceso se hará uso de las herramientas descritas más adelante en el Capítulo 3, las cuales posibilitarán la identificación de los errores cometidos en el proceso de segmentación de la arquitectura.

Por último, se pretende elaborar una prueba de concepto consistente en el despliegue del diseño final sobre una placa FPGA (Field Programmable Gate Array), de modo que pueda comprobarse sobre hardware real el buen funcionamiento y desempeño de este diseño.

%%%%%%%%%%%%%%%%%%%%% Planificación y metodología seguidos %%%%%%%%%%%%%%%%%%%%

\section{Planificación y metodología seguidos}

\subsection{Formación previa}

En primer lugar se cursará, de manera autónoma, una formación relativa a los lenguajes Scala y Chisel\footnote{Disponible en: \url{www.github.com/freechipsproject/chisel-bootcamp}}. Esta formación fue creada por diversos investigadores y colaboradores de la Universidad de California, Berkeley, y consiste en un cuaderno de Jupyter en el que se abordan, incialmente, aspectos introductorios relativos al desarrollo de software en el lenguaje Scala para, posteriormente, proporcionar al usuario nociones fundamentales sobre el diseño de hardware empleando el lenguaje Chisel en lecciones de complejidad incremental.

\subsection{Construcción del computador monociclo}

La segunda fase del proyecto consistirá en la implementación de las partes incompletas en la arquitectura de partida. Ello implicará la puesta en práctica y consiguiente asentamiento de los conocimientos adquiridos durante la realización del bootcamp sobre Scala y Chisel, además de que constituirá en sí mismo un proceso de descubrimiento de la arquitectura del núcleo monociclo y las particularidades de su implementación.

\subsection{Construcción del procesador segmentado}

A continuación, con el diseño base completado, se procederá a segmentar el núcleo en las cinco etapas siguientes:

\begin{enumerate}
  \item{Lectura de instrucción o fase de \textit{fech}}
  \vspace{-0.2cm}
  \item{Decodificación de la instrucción}
  \vspace{-0.2cm}
  \item{Ejecución}
  \vspace{-0.2cm}
  \item{Lectura/escritura de resultados en memoria}
  \vspace{-0.2cm}
  \item{Escritura en el banco de registros, o fase de \textit{writeback}}
\end{enumerate}

La integración de los registros de etapa requerirá un análisis minucioso del conexionado interno del procesador, e implicará la alternancia contínua de fases de desarrollo (construcción y adhesión de los registros al pipeline) y fases de prueba hasta verificar que la ejecución de los binarios cargados desemboca en los resultados esperados a nivel del estado arquitectónico.

\subsection{Despliegue del diseño final sobre un FPGA}

Para concluir el proyecto se elaboró una prueba de concepto consistente en el despliegue del diseño elaborado sobre una placa FPGA. Para ello se implementó, también en lenguaje Chisel, una interfaz con la que poder comunicar el computador con los dispositivos presentes en la placa. Ello ha posibilitado el envío (por medio de un módulo UART), la carga en memoria y la ejecución de binarios compilados para la arquitectura RISC-V.